\chapter{How Pay Works}\label{ch:in:how-pay-works}

Each March, New York's State Comptroller estimates the average bonus paid to the city's securities industry employees
during the bonus season just ended. For 2021 it was \$257\,500; for 2022, \$180\,000, 30\% less; for 2024, \$244\,700,
31.5\% more than the year before. An average bonus is a headline, and it hides the structure that decides what pay is
worth to the person who receives it: how much of it is fixed, how much depends on a formula or on a manager's judgement,
how much is deferred and in what, and what is lost on leaving. This chapter takes pay apart from the employee's side.
Book~16, chapter~10, designed the same machinery from the firm's side, with its bonus pools, payout formulas and
deferral schedules; here the question is what an offer is worth, and to whom.

\section{Base and bonus: the two-part structure and why it exists}

\begin{definition}[Base salary, total compensation]\label{def:in:how-pay-works:base}
\emph{Base salary}\index{base salary} is the fixed annual pay set in the employment contract and paid in instalments
whatever the year's results. \emph{Total compensation}\index{total compensation} is base salary plus the variable pay
awarded for a year, cash and deferred, at the value at which it is awarded, with benefits stated separately; it is the
figure offers are compared on, and it counts deferred awards before any \omterm{def:in:how-pay-works:rsu}{forfeiture} or change in value.
\end{definition}

Pay in trading has two parts because the firm's revenue has two parts: a cost base it must meet in any year, and a
result that swings with markets (chapter~12 measured the swing by kind of firm). A variable award lets the firm share the
good years and cut in the bad ones without cutting headcount, which is Book~16's compensation ratio at work. For the
employee the same structure transfers risk: the base is a floor, and the rest moves with the firm's year, the desk's
and the individual's.

\begin{dated}{2026-03}{New York City's securities industry bonus}\label{dat:in:how-pay-works:nyc}
\footnotesize
\begin{tabular}{@{}lrrrrr@{}}
\toprule
bonus year & 2021 & 2022 & 2023 & 2024 & 2025\\
\midrule
average bonus, first estimate (\$) & 257\,500 & 180\,000 & 176\,500 & 244\,700 & 246\,900\\
\bottomrule
\end{tabular}

\smallskip\noindent New York State Comptroller, annual estimates published each March (2022--2026); each release's own
figure for its year, except 2022, which is the figure the 2024 release gives for the previous year. The 2026 release
reports the 2025 average as 6\% above 2024's, so the 2024 estimate had been revised down. The bonus pool was a record
\$49.2 billion for 2025. In 2024 the city's average securities industry salary, bonuses included, was \$505\,677, and
bonuses were about 42\% of the industry's wages.
\end{dated}

Two things in the box matter for what follows. The average mixes every job in the industry, from operations to senior
traders, so it describes no one; and bonuses were about two-fifths of the industry's wages, so the variable part is not
a detail. Chapter~14 gives the levels; this chapter gives the structure.

\section{Formulaic and discretionary pay from the employee's side}

A variable award is either computed by a formula from a measured result, or set by a manager's judgement within a pool
(Book~16, chapter~10, defines the \emph{formulaic payout} and the \emph{discretionary bonus}). The platforms of
chapter~5 pay portfolio managers a payout rate of their book's profit (Book~16, chapter~3); a bank's or a market
maker's employee is usually paid from a discretionary pool. From the employee's side the difference is three things.
\begin{itemize}
\item \textbf{Shape.} A payout of a share of positive profit and nothing below zero is an option on the book's result:
convex, with a long right tail and a mass at zero. A discretionary award is smoother, because the manager spreads the
pool and weighs more than one year.
\item \textbf{Whose risk.} A formula exposes the employee to the book's result, including luck; discretion exposes the
employee to the pool, which follows the firm's and the desk's year, and to the manager's judgement.
\item \textbf{What ends it.} A formulaic book can be closed by a loss limit (chapter~5's drawdown ladders), ending the
job and the payout together; a discretionary employee's award can go to zero in a bad year without the job ending.
\end{itemize}
The certainty equivalent (Book~4, chapter~9) turns these differences into money: the sure amount a person with a
stated relative risk aversion would accept in place of the uncertain package. Two people offered the same package can
value it differently, because their risk aversion and their other wealth differ.

\section{Deferral, vesting, forfeiture and clawback: what deferred pay is worth}

Part of a variable award is often deferred: awarded now, paid later in tranches under a vesting schedule, and held
meanwhile in cash, in the firm's shares or in units of its funds (Book~16, chapter~10, defines deferred compensation,
the vesting schedule, malus and clawback).

\begin{definition}[Restricted stock unit, forfeiture, good-leaver provision]\label{def:in:how-pay-works:rsu}
A \emph{restricted stock unit}\index{restricted stock unit} (RSU) is a promise to deliver one share of the employer's
stock, or its cash value, when the unit vests, subject to conditions such as continued employment. \emph{Forfeiture}
\index{forfeiture} is the loss of unvested awards, usually on leaving the firm voluntarily or on dismissal for cause.
A \emph{good-leaver provision}\index{good-leaver provision} is a term of an award that lets some leavers, typically on
retirement, death, disability or redundancy, keep unvested awards, which then vest on the original schedule or at once.
\end{definition}

\begin{dated}{2025-10}{Deferral rules for banks and one bank's awards}\label{dat:in:how-pay-works:rules}
\textbf{EU} (Directive 2013/36/EU, Article~94, as amended in 2019): for staff whose work has a material impact on the
firm's risk, variable pay may not exceed 100\% of fixed pay (200\% with shareholder approval); at least 50\% of variable
pay is in shares or equivalent instruments; at least 40\% is deferred over not less than four to five years, vesting no
faster than pro rata, and at least 60\% of a particularly high amount. \textbf{UK} (PRA PS21/25 and FCA PS25/15,
in force for performance years starting after 16 October 2025): a four-year minimum deferral for all material risk
takers; 40\% deferral on the first \pounds660\,000 of a bonus and 60\% on the part above it; the bonus cap was removed
from 31 October 2023. \textbf{One US bank} (its Form 10-K for 2025): RSUs without performance conditions ``generally vest
and underlying shares of common stock are delivered (net of required withholding tax) over a three-year period'', with
vesting accelerated on ``retirement, death, disability and, in certain cases, conflicted employment''.
\end{dated}

What a deferred award is worth to the employee depends on four things, which the chapter's build values together.
\begin{enumerate}
\item \textbf{Time.} A dollar in four years is worth less than a dollar now; at 5\% a year, a tranche vesting in four
years is worth 82 cents of award.
\item \textbf{The instrument.} Deferred shares carry the stock's risk: with 30\% volatility, a tranche's value at
vesting ranges widely even with no expected change.
\item \textbf{Leaving.} Unvested tranches are forfeited on a voluntary move, so their value is weighted by the
probability of staying; a \omterm{def:in:how-pay-works:rsu}{good-leaver provision} removes that weight for the leavers it covers.
\item \textbf{Malus and clawback.} An award can be reduced before vesting (malus) or recovered after it (clawback) for
misconduct or a later loss; the chapter's model treats these as rare and states them rather than pricing them.
\end{enumerate}
Tax follows the vesting, not the award. In the United States, property subject to ``a substantial risk of \omterm{def:in:how-pay-works:rsu}{forfeiture}''
is not income ``until it becomes substantially vested'', and then its market value is; a deferred award's tax is paid
on its value when it vests, which is why a bank delivers shares ``net of required withholding tax''. Chapter~15 gives the
rules of other locations.

\begin{example}[The handcuff]\label{ex:in:how-pay-works:handcuff}
The chapter's illustrative bank offer defers 40\% of each award into the firm's stock, vesting in equal tranches over
four years. For someone who stays, the unvested balance at grant value grows to about \$90\,000 at the end of the first
year, \$158\,000 after two, \$203\,000 after three and levels off near \$226\,000 from the fourth, when each year's new
deferral is matched by vesting (\cref{fig:in:how-pay-works:unvested}). Leaving at the end of year two forfeits a present
value of about \$129\,000; at the end of year four, \$170\,000.
\end{example}

\begin{omfigure}
\begin{tikzpicture}
\begin{axis}[width=10.5cm,height=5.4cm,ybar,area legend,bar width=10pt,xmin=0.4,xmax=5.6,ymin=0,ymax=260,xtick={1,2,3,4,5},
  xlabel={\small end of year},ylabel={\small \$ thousand},tick label style={font=\footnotesize},grid=major,
  grid style={gray!20},table/col sep=comma,legend style={legend cell align=left,font=\scriptsize,at={(0.5,-0.3)},anchor=north,
  /tikz/every even column/.append style={column sep=8pt}},legend columns=2]
\addplot[fill=omDef!55,draw=omDef] table[x=year,y=unvested]{figdata/industry/13-how-pay-works/bank_unvested.csv};
\addplot[fill=omThm!55,draw=omThm] table[x=year,y=forfeit]{figdata/industry/13-how-pay-works/bank_unvested.csv};
\legend{unvested balance at grant value (stayer),present value forfeited by leaving then}
\end{axis}
\end{tikzpicture}

\omcaption{The illustrative bank offer's deferred awards: 40\% of each year's award deferred into stock and vesting over
four years. The balance a leaver walks away from builds for four years, then levels off. Parameters are illustrative;
data: 20\,000 simulated careers, through \texttt{in\_pay.bank\_unvested} and \texttt{forfeit\_by\_year}.}\label{fig:in:how-pay-works:unvested}
\end{omfigure}

\section{Sign-ons, guarantees and buyouts}

\begin{definition}[Sign-on bonus, guaranteed bonus]\label{def:in:how-pay-works:signon}
A \emph{sign-on bonus}\index{sign-on bonus} is a payment on joining, usually repayable in full or in part if the
employee leaves within a stated period. A \emph{guaranteed bonus}\index{guaranteed bonus} is a variable award whose
minimum for a stated year, usually the first, is fixed in the offer.
\end{definition}

A hiring firm uses the two to solve two problems. The sign-on pays for what the candidate gives up by moving, and its
repayment clause keeps them for its term; the guarantee removes the first year's risk, when a new joiner has no track
record at the firm and a partial year to build one. A third instrument, the \emph{deferral buyout} of Book~16,
chapter~10, replaces the unvested awards the candidate forfeits by leaving, usually as new awards on the old vesting
dates. From the candidate's side the arithmetic is the handcuff of the example above run backwards: a candidate who
leaves the illustrative bank at the end of year two forfeits awards with a grant value of about \$158\,000, and a buyout
of that size on the same dates makes the move neutral for the deferred part. The candidate should also ask what
repayment and \omterm{def:in:how-pay-works:rsu}{forfeiture} terms apply to the buyout itself, because a buyout is usually deferred again.

The EU's rules for banks limit what a guarantee may do: guaranteed variable pay ``is exceptional, occurs only when
hiring new staff and where the institution has a sound and strong capital base and is limited to the first year of
employment''; outside the firms such rules cover it is a matter of contract. Chapter~28 returns to buyouts and notice
periods when it follows moves between firms, and Book~18 to how offers are negotiated.

\section{Partnership, equity and carried interest}

\begin{definition}[Partnership share, carried interest]\label{def:in:how-pay-works:carry}
A \emph{partnership share}\index{partnership share} is a member's entitlement to a share of a partnership's profit,
fixed by the partnership agreement, in return for capital contributed and work; it is paid as \omterm{def:in:reading-a-trading-firms-accounts:members}{members' remuneration}
(chapter~11) rather than as salary. \emph{Carried interest}\index{carried interest} is a share of a fund's profits,
typically above a hurdle, paid to its managers through an interest in the fund's general partner rather than as a fee.
\end{definition}

Ownership pays differently from employment. A partner in a limited liability partnership (Book~16, chapter~2) shares
the profit after costs, including the pay of the employees, and bears its losses up to the capital put in; chapter~11
showed that in some trading partnerships most of the profit goes to a corporate member rather than to individuals.
Equity in a private firm is worth what the firm will pay for it on the terms of its agreements, which usually restrict
sale and set a price on leaving. \omterm{def:in:how-pay-works:carry}{Carried interest} is a performance fee's share paid to the people who manage the fund
(Book~1, chapter~3, defines the fee and the high-water mark).

\begin{dated}{2026-04}{Carried interest tax}\label{dat:in:how-pay-works:carry}
\textbf{UK}: from 6 April 2026, \omterm{def:in:how-pay-works:carry}{carried interest} is taxed wholly within income tax, ``treated as trading profits and
subject to Income Tax and Class 4 National Insurance contributions''; for qualifying \omterm{def:in:how-pay-works:carry}{carried interest} ``the amount to be
treated as trading profits is 72.5\% of the qualifying profits''. \textbf{US}: gains on an applicable partnership
interest are long-term only after three years' holding rather than one (26 U.S.C.\ 1061).
\end{dated}

\section{Structure by firm type}

The same parts combine differently by kind of employer, for reasons that the earlier chapters give.
\begin{itemize}
\item \textbf{Banks.} Variable pay above rule thresholds is deferred and partly in shares, under the rules of the dated
box for the staff they cover; a US bank's own RSUs vest over three years. Deferral retains, and a buyout is
part of the cost of hiring an experienced banker (chapter~7's divisions, chapter~18's roles).
\item \textbf{Market makers and proprietary firms.} Private firms publish no pay rules; an offer is typically a base
and a discretionary award from a pool that follows the firm's revenue, which chapter~12 found moves with volatility more than
any other kind of firm's. Any deferral is contractual, not regulatory, outside the entities that the rules cover.
\item \textbf{Platforms.} A portfolio manager's pay is a payout rate of the book's profit (chapter~5); an analyst's, in
the illustrative offer of this chapter, is a base and a small share of the book's result; the book can be closed by a
loss limit. Sign-ons and guarantees buy people out of their previous firm.
\item \textbf{Asset managers and funds.} Base and bonus, which may be partly deferred into the firm's funds or shares;
senior people may share in performance fees or \omterm{def:in:how-pay-works:carry}{carried interest} (chapter~8).
\item \textbf{Exchanges, vendors and regulators.} Listed exchanges and vendors pay base and bonus like other listed
companies; a regulator pays on a published scale (chapter~9).
\item \textbf{Crypto firms.} Tokens and \omterm{def:in:crypto-native-firms:tokens}{token warrants} add a volatile, often locked-up component (chapter~10).
\end{itemize}

\section{Tutorial: three offers, one choice}\label{tut:in:how-pay-works:tut}

\begin{tutorial}
\textbf{Goal.} Value three stylised offers from the employee's side and compare them at two levels of risk aversion.
\textbf{End state:} \cref{fig:in:how-pay-works:offers,fig:in:how-pay-works:ce} and the table below.
\begin{tutsteps}
\item \textbf{The offers} (every number illustrative, in US dollars): a \emph{bank} role with base 250\,000 and a
discretionary award of median 200\,000 (log standard deviation 0.5), 40\% deferred into stock (volatility 30\%) over four
years, forfeited on leaving; a \emph{market-maker} role with base 200\,000 and an award of median 250\,000 (log standard
deviation 0.8), paid in cash; a \emph{platform analyst} with base 175\,000, a payout of 2\% of the book's profit (mean
15 million, standard deviation 20 million a year), the book closed after a year losing more than 10 million, a sign-on of
100\,000 repayable within a year, and a first-year guarantee of 150\,000.
\item \textbf{The career.} Five years, a 10\% chance a year of leaving voluntarily, a 5\% discount rate;
\texttt{firm.payoffer.simulate} draws 20\,000 careers per offer (\cref{lst:in:how-pay-works:sim}).
\omcode{code/firm/payoffer/firm_payoffer.py}{87}{103}{Each year: base while working, the award's upfront part, and each
deferred tranche delivered if still employed or a good leaver, forfeited otherwise.}\label{lst:in:how-pay-works:sim}
\item \textbf{Value.} \texttt{summary} gives the mean and percentiles of the present value received;
\texttt{certainty\_equivalent(pv, rra, 500\_000)} gives the certainty equivalent with 500\,000 of other wealth.
\end{tutsteps}
\end{tutorial}

\begin{center}\footnotesize
\begin{tabular}{@{}lrrrrr@{}}
\toprule
offer (illustrative) & mean PV & 10th pct & 90th pct & CE, RRA 1 & CE, RRA 3\\
\midrule
bank & 1\,618 & 597 & 2\,230 & 1\,500 & 1\,186\\
market maker & 1\,949 & 686 & 2\,956 & 1\,779 & 1\,365\\
platform analyst & 1\,703 & 171 & 3\,010 & 1\,439 & 904\\
\bottomrule
\end{tabular}

\smallskip Present value of five years' pay, \$ thousand; CE is the certainty equivalent.
\end{center}

\begin{omfigure}
\begin{tikzpicture}
\begin{axis}[width=10.5cm,height=5.4cm,xmin=0.4,xmax=3.6,ymin=0,ymax=3400,xtick={1,2,3},
  xticklabels={bank,market maker,platform analyst},ylabel={\small present value, \$ thousand},
  tick label style={font=\footnotesize},grid=major,grid style={gray!20},table/col sep=comma,
  legend style={legend cell align=left,font=\scriptsize,at={(0.5,-0.2)},anchor=north,/tikz/every even column/.append style={column sep=8pt}},
  legend columns=2]
\addplot[draw=none,mark=none,forget plot,error bars/.cd,y dir=both,y explicit,error mark=none,
  error bar style={omDef,line width=1pt}]
  table[x=pos,y=p50,y error plus expr=\thisrow{p90}-\thisrow{p50},y error minus expr=\thisrow{p50}-\thisrow{p10}]{figdata/industry/13-how-pay-works/dist.csv};
\addplot[draw=none,mark=none,forget plot,error bars/.cd,y dir=both,y explicit,error mark=none,
  error bar style={omDef!35,line width=14pt}]
  table[x=pos,y=p50,y error plus expr=\thisrow{p75}-\thisrow{p50},y error minus expr=\thisrow{p50}-\thisrow{p25}]{figdata/industry/13-how-pay-works/dist.csv};
\addplot[only marks,mark=-,mark size=7pt,line width=1.4pt,omDef] table[x=pos,y=p50]{figdata/industry/13-how-pay-works/dist.csv};
\addplot[only marks,mark=diamond*,mark size=2.5pt,omThm] table[x=pos,y=mean]{figdata/industry/13-how-pay-works/dist.csv};
\legend{median,mean}
\end{axis}
\end{tikzpicture}

\omcaption{Five years of pay under three illustrative offers: boxes span the 25th to 75th percentiles of the present value
received, whiskers the 10th to 90th, the bar is the median and the diamond the mean. The platform analyst's spread is the widest: a closed book
ends the job and the payout together. Parameters are illustrative; data: \texttt{in\_pay.table}.}\label{fig:in:how-pay-works:offers}
\end{omfigure}

\begin{omfigure}
\begin{tikzpicture}
\begin{axis}[width=10.5cm,height=5.4cm,xmin=0.4,xmax=5.1,ymin=500,ymax=2000,xlabel={\small relative risk aversion},
  ylabel={\small certainty equivalent, \$ thousand},tick label style={font=\footnotesize},grid=major,grid style={gray!20},
  table/col sep=comma,legend style={legend cell align=left,font=\scriptsize,at={(0.97,0.97)},anchor=north east}]
\addplot[omDef,thick,mark=*] table[x=rra,y=bank]{figdata/industry/13-how-pay-works/ce_curve.csv};
\addplot[omThm,thick,mark=square*] table[x=rra,y=maker]{figdata/industry/13-how-pay-works/ce_curve.csv};
\addplot[black!70,thick,mark=triangle*] table[x=rra,y=platform]{figdata/industry/13-how-pay-works/ce_curve.csv};
\legend{bank,market maker,platform analyst}
\end{axis}
\end{tikzpicture}

\omcaption{Certainty equivalents of the three illustrative offers against relative risk aversion, with 500\,000 of other
wealth. The platform package loses most as risk aversion rises: at low aversion it is worth about as much as the bank
offer, at high aversion much less. Data: \texttt{in\_pay.ce\_curve}.}\label{fig:in:how-pay-works:ce}
\end{omfigure}

The market-maker offer is worth most at every level of risk aversion here, because its parameters give it the highest
median award; that is a property of the illustrative numbers, not of market makers. The lesson is in the shapes. The
platform package's mean is above the bank's, but its 10th percentile is \$171\,000, because a book closed after a bad year
ends the pay; in 35.4\% of simulated careers the book is closed within five years. Its certainty equivalent falls from
\$1\,439\,000 at relative risk aversion~1 to \$904\,000 at~3, while the bank's falls from \$1\,500\,000 to \$1\,186\,000.
A candidate who cannot bear the platform's lower tail should value it well below its mean.

\textbf{What to change next.} Give the bank offer a \omterm{def:in:how-pay-works:rsu}{good-leaver provision} and see how much of the \omterm{def:in:how-pay-works:rsu}{forfeiture} it removes
for a planned move; raise the platform's loss limit and trace the certainty equivalent.

\section{Build: the offer valuer}\label{bld:in:how-pay-works:payoffer}

\begin{build}
\bfield{Purpose} Value a pay package from the employee's side, for this chapter and for chapters~22, 28 and~30 (and for
Book~18's chapter on offers).
\bfield{Interface} \texttt{firm.payoffer}: \texttt{Package(name, base, median, sigma, deferral, vest\_years, inst\_vol,
good\_leaver, sign\_on, sign\_on\_years, guarantee, share, pnl\_mean, pnl\_sd, cut)};
\texttt{simulate(pkg, years, hazard, rate, n, rng, leave\_year)}; \texttt{summary(sim)};
\texttt{certainty\_equivalent(pv, rra, wealth)}; \texttt{buyout(pkg, year, n, rng, rate)}.
\bfield{Rules} Awards are paid at the end of the year worked; deferred tranches vest in equal parts after the award and
are delivered only to those still employed or covered as good leavers; a sign-on is repaid in full within its period;
a guarantee floors the first award; a closed book pays nothing for its last year; every parameter is the caller's.
\bfield{Acceptance tests} \texttt{code/firm/payoffer/tests/}: a base-only package is an annuity; deferral without risk
matches hand arithmetic; a leaver forfeits exactly the unvested tranches and a good leaver none; a sign-on is repaid on
an early exit; the buyout equals the unvested balance; the certainty equivalent of a sure amount is that amount and falls
with risk aversion; the mean with stock risk matches the closed form within 1\%.
\bfield{Stretch} Malus and clawback as events with a probability; taxes by location (chapter~15); correlated awards
across years; a leaving hazard that depends on the year's award.
\end{build}

\begin{omsources}
\item Office of the New York State Comptroller, securities industry bonus releases, March 2022, 2024, 2025 and 2026.
\item Directive 2013/36/EU, Article~94, and Directive (EU) 2019/878; Bank of England, PS21/25 -- Remuneration Reform.
\item The Goldman Sachs Group, Inc., Form 10-K for 2025, Note~29.
\item US Internal Revenue Service, Publication 525; 26 U.S.C.\ 1061.
\item HM Treasury and HMRC, Reform of the tax treatment of carried interest: revised tax regime (2025).
\end{omsources}

\section{Exercises}

\begin{exercise}[$\star$]\label{exo:in:how-pay-works:1}
By how much did New York City's average securities bonus change from 2021 to 2022 and from 2023 to 2024?
\end{exercise}

\begin{exercise}[$\star$]\label{exo:in:how-pay-works:2}
Under the UK rules in force for performance years after October 2025, how much of a \pounds1\,000\,000 bonus must be
deferred?
\end{exercise}

\begin{exercise}[$\star$]\label{exo:in:how-pay-works:3}
At a 5\% discount rate, what is a tranche vesting in four years worth per unit of award, before any risk?
\end{exercise}

\begin{exercise}[$\star\star$]\label{exo:in:how-pay-works:4}
In the bank offer, list the tranches still unvested at the end of year two and compute their grant value, using an
expected award of $200\,000\,e^{0.5^2/2}$.
\end{exercise}

\begin{exercise}[$\star\star$]\label{exo:in:how-pay-works:5}
Why is a payout of a share of positive profit an option, and what does that do to its certainty equivalent?
\end{exercise}

\begin{exercise}[$\star\star$]\label{exo:in:how-pay-works:6}
A candidate is offered a sign-on of \$100\,000 repayable in full if she leaves within a year. She thinks there is a 15\%
chance she will. What is the sign-on worth to her in expectation, ignoring discounting?
\end{exercise}

\begin{exercise}[$\star\star\star$]\label{exo:in:how-pay-works:7}
\emph{Coding.} Give the bank offer a \omterm{def:in:how-pay-works:rsu}{good-leaver provision} covering voluntary moves. How much does the expected
\omterm{def:in:how-pay-works:rsu}{forfeiture} fall, and how much does the mean present value rise?
\end{exercise}

\begin{exercise}[$\star\star\star$]\label{exo:in:how-pay-works:8}
\emph{Find the flaw.} ``The average Wall Street bonus was \$246\,900 in 2025, so a trader joining a bank in New York
should expect about that.''
\end{exercise}

\section{Problem: Three Offers, One Choice}

\begin{problem}[{Weekend problem --- three offers, one choice}]\label{pb:in:how-pay-works:1}
A quant with 500\,000 of savings holds three offers: a bank, a market maker and a platform. She wants to know what each
is worth to her, and what the bank would cost her to leave in two years.

\medskip\noindent\textbf{Part I --- The parts.}
\begin{enumerate}
\item Define \omterm{def:in:how-pay-works:base}{base salary} and \omterm{def:in:how-pay-works:base}{total compensation}. Why do offers compare on \omterm{def:in:how-pay-works:base}{total compensation}, and what does it omit?
\item Give New York City's average securities bonus for 2021--2025 and the share of wages bonuses made up in 2024.
\item Why does pay in trading have two parts?
\item Contrast formulaic and discretionary pay in shape, whose risk, and what ends it.
\item Define the certainty equivalent and say what it depends on besides the package.
\end{enumerate}

\medskip\noindent\textbf{Part II --- Deferral.}
\begin{enumerate}[resume]
\item Define an RSU, \omterm{def:in:how-pay-works:rsu}{forfeiture} and a \omterm{def:in:how-pay-works:rsu}{good-leaver provision}.
\item State the EU and UK deferral rules and one US bank's RSU vesting.
\item Name four things that set a deferred award's value to the employee.
\item When is a US employee's RSU taxed, and on what value?
\item Give the bank offer's unvested balance at the end of each year and the value forfeited by leaving after year two.
\end{enumerate}

\medskip\noindent\textbf{Part III --- The offers.}
\begin{enumerate}[resume]
\item Define a sign-on and a \omterm{def:in:how-pay-works:signon}{guaranteed bonus}, and say what each solves for the firm.
\item State the three offers' parameters and the career assumptions.
\item Give each offer's mean present value and 10th and 90th percentiles.
\item In what share of careers is the platform book closed within five years?
\item Define \omterm{def:in:how-pay-works:carry}{carried interest} and a \omterm{def:in:how-pay-works:carry}{partnership share}, and state the UK and US carried-interest tax rules.
\end{enumerate}

\medskip\noindent\textbf{Part IV --- The verdict.}
\begin{enumerate}[resume]
\item State the \emph{named result}: each offer's certainty equivalent at relative risk aversion 1 and 3, and the
deferral buyout that makes moving from the bank in year two neutral.
\item Why does the platform offer lose most as risk aversion rises?
\item What should she ask about a buyout besides its size?
\item Which result depends on the illustrative parameters, and which on the structure?
\item In two sentences, how should she choose?
\end{enumerate}
\end{problem}

\section{Interview questions}

\begin{interviewq}[$\star$ \iqroles{trader, researcher}]\label{iq:in:how-pay-works:1}
Your bonus is 2\% of your book's positive P\&L. What option is that, and on what?
\end{interviewq}

\begin{interviewq}[$\star$ \iqroles{bank}]\label{iq:in:how-pay-works:2}
Why do bank regulators require part of variable pay to be deferred and paid in shares?
\end{interviewq}

\begin{interviewq}[$\star\star$ \iqroles{researcher, risk}]\label{iq:in:how-pay-works:3}
Compute the certainty equivalent, with log utility and wealth 100, of a coin flip paying 0 or 200.
\end{interviewq}

\begin{interviewq}[$\star\star$ \iqroles{developer}]\label{iq:in:how-pay-works:4}
Design the data model for an employee's awards so that the unvested balance on any date, and the \omterm{def:in:how-pay-works:rsu}{forfeiture} on any
leaving date, can be computed.
\end{interviewq}

\begin{interviewq}[$\star\star$ \iqroles{bank, risk}]\label{iq:in:how-pay-works:5}
What is the difference between malus and clawback, and which is harder to enforce?
\end{interviewq}

\begin{interviewq}[$\star\star\star$ \iqroles{researcher}]\label{iq:in:how-pay-works:6}
A firm defers 40\% of pay over four years. Estimate how much that saves it in turnover, and what data you would need.
\end{interviewq}
